\begin{figure}[ht]
		\centering
		\begin{tikzpicture}[>=Stealth, line width=0.8pt]
			
			\def\L{3.6}
			\def\dy{1.50}
			
			% ============================================================
			% SP1 — cerniera in A + carrello in C
			%        (rimossi: momento in A, cerniera in B)
			% ============================================================
			
			\begin{scope}[yshift=0cm]
				
				
				\draw[line width=1.5pt] (0,0) -- (\L,0);
				
				
				% incastro in A rimosso (grigio) — disegnato per primo
				\draw[gray!50, thick] (0,-0.38) -- (0, 0.38);
				\foreach \y in {-0.28,-0.14,0,0.14,0.28}{
					\draw[gray!50, thin] (0,\y) -- (-0.20,\y-0.13);
				}
				
				% cerniera in A (attiva) — triangolo
				\draw[thick] (0,0) -- (-0.20,-0.34)
				-- (0.20,-0.34) -- cycle;
				\draw[thick] (-0.30,-0.34) -- (0.30,-0.34);
				\foreach \x in {-0.22,-0.11,0,0.11,0.22}{
					\draw[thin] (\x,-0.34) -- (\x-0.10,-0.50);
				}
				
				% carrello verticale in C (attivo)
				\draw[thick] (\L/2,0) -- (\L/2-0.20,-0.34)
				-- (\L/2+0.20,-0.34) -- cycle;
				\draw[thick] (\L/2-0.14,-0.40) circle (0.07);
				\draw[thick] (\L/2+0.14,-0.40) circle (0.07);
				\draw[thick] (\L/2-0.30,-0.50) -- (\L/2+0.30,-0.50);
				\foreach \x in {-0.22,-0.11,0,0.11,0.22}{
					\draw[thin] (\L/2+\x,-0.50)
					-- (\L/2+\x-0.10,-0.64);
				}
				% mezzaluna sopra il vertice del triangolo in C
				\fill[white]      (\L/2-0.08,0) arc(180:360:0.08) -- cycle;
				\draw[thick]      (\L/2-0.08,0) arc(180:360:0.08);
				
				% cerniera in B rimossa (grigio)
				\draw[gray!50, thick] (\L,0) -- (\L-0.20,-0.34)
				-- (\L+0.20,-0.34) -- cycle;
				\draw[gray!50, thick] (\L-0.30,-0.34)
				-- (\L+0.30,-0.34);
				\foreach \x in {-0.22,-0.11,0,0.11,0.22}{
					\draw[gray!50, thin] (\L+\x,-0.34)
					-- (\L+\x-0.10,-0.50);
				}
				
				% pallini A e B — disegnati per ultimi, sopra i vincoli
				\draw[fill=white, thick]       (0,0)  circle (2pt);
				\draw[gray!50, fill=white, thick] (\L,0) circle (2pt);
				
				% etichette dei punti
				\node[above right,  font=\scriptsize]            at (0,0)    {$A$};
				\node[above,        font=\scriptsize]            at (\L/2,0) {$C$};
				\node[above, font=\scriptsize, gray!60]    at (\L,0)   {$B$};
				
				% etichetta del sistema principale e nomenclatura iperstatiche
				\node[font=\scriptsize, anchor=east] at (-0.5,0) {SP1};
				\node[font=\scriptsize, anchor=west] at (\L+0.3,0)
				{$\chi_1\!\equiv\!\mu_A,\;
					\chi_2\!\equiv\!X_B,\;
					\chi_3\!\equiv\!Y_B$};
				
			\end{scope}
			
			% ============================================================
			% SP2 — mensola (incastro in A)
			%        (rimossi: carrello in C, cerniera in B)
			% ============================================================
			\begin{scope}[yshift=-\dy cm]
				
				\draw[line width=1.5pt] (0,0) -- (\L,0);
				
				% incastro in A (attivo) — niente perno: l'incastro non ha pallino
				\draw[thick] (0,-0.38) -- (0, 0.38);
				\foreach \y in {-0.28,-0.14,0,0.14,0.28}{
					\draw[thin] (0,\y) -- (-0.20,\y-0.13);
				}
				
				% carrello verticale in C rimosso (grigio)
				\draw[gray!50, thick] (\L/2,0) -- (\L/2-0.20,-0.34)
				-- (\L/2+0.20,-0.34) -- cycle;
				\draw[gray!50, thick] (\L/2-0.14,-0.40) circle (0.07);
				\draw[gray!50, thick] (\L/2+0.14,-0.40) circle (0.07);
				\draw[gray!50, thick] (\L/2-0.30,-0.50)
				-- (\L/2+0.30,-0.50);
				\foreach \x in {-0.22,-0.11,0,0.11,0.22}{
					\draw[gray!50, thin] (\L/2+\x,-0.50)
					-- (\L/2+\x-0.10,-0.64);
				}
				% mezzaluna grigia sopra il vertice del triangolo in C (vincolo rimosso)
				\fill[white]            (\L/2-0.08,0) arc(180:360:0.08) -- cycle;
				\draw[gray!50, thick]   (\L/2-0.08,0) arc(180:360:0.08);
				
				% cerniera in B rimossa (grigio)
				\draw[gray!50, thick] (\L,0) -- (\L-0.20,-0.34)
				-- (\L+0.20,-0.34) -- cycle;
				\draw[gray!50, thick] (\L-0.30,-0.34)
				-- (\L+0.30,-0.34);
				\foreach \x in {-0.22,-0.11,0,0.11,0.22}{
					\draw[gray!50, thin] (\L+\x,-0.34)
					-- (\L+\x-0.10,-0.50);
				}
				
				% pallino in B — disegnato per ultimo, sopra il triangolo grigio
				\draw[gray!50, fill=white, thick] (\L,0) circle (2pt);
				
				% etichette dei punti
				\node[above right,  font=\scriptsize]              at (0,0)    {$A$};
				\node[above,        font=\scriptsize, gray!60]    at (\L/2,0) {$C$};
				\node[above, font=\scriptsize, gray!60]     at (\L,0)   {$B$};
				
				% etichetta del sistema principale e nomenclatura iperstatiche
				\node[font=\scriptsize, anchor=east] at (-0.5,0) {SP2};
				\node[font=\scriptsize, anchor=west] at (\L+0.3,0)
				{$\chi_1\!\equiv\!X_B,\;
					\chi_2\!\equiv\!Y_B,\;
					\chi_3\!\equiv\!Y_C$};
				
			\end{scope}
			
			% ============================================================
			% SP3 — cerniera in B + carrello in C
			%        (rimosso: incastro in A)
			% ============================================================
			\begin{scope}[yshift=-2*\dy cm]
				
				\draw[line width=1.5pt] (0,0) -- (\L,0);
				
				% incastro in A rimosso (grigio)
				\draw[gray!50, thick] (0,-0.38) -- (0, 0.38);
				\foreach \y in {-0.28,-0.14,0,0.14,0.28}{
					\draw[gray!50, thin] (0,\y) -- (-0.20,\y-0.13);
				}
				
				% carrello verticale in C (attivo)
				\draw[thick] (\L/2,0) -- (\L/2-0.20,-0.34)
				-- (\L/2+0.20,-0.34) -- cycle;
				\draw[thick] (\L/2-0.14,-0.40) circle (0.07);
				\draw[thick] (\L/2+0.14,-0.40) circle (0.07);
				\draw[thick] (\L/2-0.30,-0.50) -- (\L/2+0.30,-0.50);
				\foreach \x in {-0.22,-0.11,0,0.11,0.22}{
					\draw[thin] (\L/2+\x,-0.50)
					-- (\L/2+\x-0.10,-0.64);
				}
				% mezzaluna nera sopra il vertice del triangolo in C (vincolo attivo)
				\fill[white]      (\L/2-0.08,0) arc(180:360:0.08) -- cycle;
				\draw[thick]      (\L/2-0.08,0) arc(180:360:0.08);
				
				% cerniera in B (attiva)
				\draw[thick] (\L,0) -- (\L-0.20,-0.34)
				-- (\L+0.20,-0.34) -- cycle;
				\draw[thick] (\L-0.30,-0.34) -- (\L+0.30,-0.34);
				\foreach \x in {-0.22,-0.11,0,0.11,0.22}{
					\draw[thin] (\L+\x,-0.34) -- (\L+\x-0.10,-0.50);
				}
				
				% pallino in B — disegnato per ultimo, sopra il triangolo nero
				\draw[fill=white, thick] (\L,0) circle (2pt);
				
				% etichette dei punti
				\node[above right,  font=\scriptsize, gray!60]   at (0,0)    {$A$};
				\node[above,        font=\scriptsize]            at (\L/2,0) {$C$};
				\node[above, font=\scriptsize]             at (\L,0)   {$B$};
				
				% etichetta del sistema principale e nomenclatura iperstatiche
				\node[font=\scriptsize, anchor=east] at (-0.5,0) {SP3};
				\node[font=\scriptsize, anchor=west] at (\L+0.3,0)
				{$\chi_1\!\equiv\!X_A,\;
					\chi_2\!\equiv\!Y_A,\;
					\chi_3\!\equiv\!\mu_A$};
				
			\end{scope}
			
		\end{tikzpicture}
		\caption{Trave iperstatica di grado 3
			(\ESEref{ex:iper_stat_3}): tre sistemi principali
			isostatici. I vincoli rimossi sono in grigio;
			le corrispondenti reazioni diventano le incognite
			iperstatiche $\chi_1$, $\chi_2$, $\chi_3$.}
		\label{fig:iper_stat_3}
	\end{figure}
